We study the generalized star height of terms of Kleene algebra with complement, interpreted over $\RSUB$, the finite relational models in which variables denote forward relations on finite linear orders.
We ask how simple the underlying structures can be while height one and height two are still separated.
We measure simplicity by the pathwidth of the Gaifman graph of a valuation together with the successor relation of the order.
Pathwidth one essentially means words, and there the question is closely related to the open generalized star-height problem for regular languages.
We show that pathwidth two already suffices for a separation.
We give a fixed term $\Psi$ of star height two over three variables that is not equivalent to any term of star height at most one on the class of $\RSUB$ valuations of pathwidth at most two.
The counter-models are very simple: their Gaifman graphs are paths with a few fans attached.
On pathwidth one, $\Psi$ itself is equivalent to a term of star height one, so the bound two is tight for $\Psi$.
The proof combines Ehrenfeucht--Fra\"iss\'e games with an aperiodic and a non-aperiodic finite monoid.
These give, for every depth, two modules that no star of that depth can tell apart, while a single nested star reads a parity bit from them.
A replacement lemma and a parity argument for first-order logic with modular predicates then rule out all terms of star height one.
